\section{Zadanie 11}
Udowodnij, że jeśli \(\Var X = 0\), to istnieje \(a\)
takie, że \(P(X=a)=1\)

\subsection*{Rozwiązanie}
Skorzystamy z zadania 4. z pierwszego zestawu - dla ciągu zdarzeń
\(\{A_n\}\), jeśli \(\forall_{i\in\mathbb{N}} P(A_i) = 1\), to
\(P(\bigcap_{i\in\mathbb{N}} A_i) = 1\)

Niech \(a = E[X]\) (wartość oczekiwana istnieje, bo wariancja istnieje)

\subsubsection*{Lemat}
\[\forall \varepsilon > 0: P(|X-a| > \varepsilon) = 0\]

Zauważmy:
\[\Var X \geqslant
	\varepsilon^2 \cdot P(|X-a| > \varepsilon)\]
Wariancja jest sumą (całką) kwadratów odchyleń przyjmowanych wartości
od wartości oczekiwanej. Uwzględniając tylko te elementy, które nie
należą do zbioru \([a-\varepsilon, a+\varepsilon]\), dostajemy dolne
ograniczenie na wariancję.

Wiedząc że wariancja wynosi 0, to prawdopodobieństwo także musi być
zerowe. \(\quad\square\)

Przekształcając:
\[\forall \varepsilon > 0: P(|X-a| \leqslant \varepsilon) = 1\]
\[\forall \varepsilon > 0: P(X \in [a-\varepsilon, a+\varepsilon]) = 1\]

Podstawiając dla
\(\varepsilon = 1, \frac1{2}, \frac1{3}, \frac1{4}, \ldots\)
\[\forall_{n\in\mathbb{N_+}}
	P\left(X\in\left[a-\frac1{n}, a+\frac1{n}\right]\right) = 1\]
Teraz korzystając z zadania 1/4:
\[P(X\in\{a\}) = 1\quad\square\]
